\section{The data specification}
\label{sec:spec}
%==========================================================================

\emph{Every row is generated from the model code and is a consequence of something proved
above. It has two tiers, and the order is not cosmetic: Tier~1 is a single measurement
nobody has published, and until it reports there is no defensible answer to the question
Tier~2 asks.}

\begin{adjustwidth}{-1.4cm}{-1.4cm}
\begin{table}[htbp]
\centering
\footnotesize
\caption{\label{tab:spec}What a record would have to contain for $\theta$ to be
identifiable. Quantities are in observable units, not multiples of $K_m$, because $K_m$ is
not known --- and by the last row does not need to be.}
\begin{tabular}{@{}lll@{}}
\toprule
\textbf{quantity} & \textbf{requirement} & \textbf{why} \\
\midrule
marker & plasma 3-hydroxypropionate & endpoint association (RR 0.7, 95\% CI 0.6--0.9) and immunity to the carnitine sign reversal that disqualifies C3 \\
second channel & $n$-propionylglycine, optional & a glycine conjugate, so carnitine does not raise it. Worth adding only after the load record is good enough for the assay term to bind \\
sampling interval & $\le$ 4 days through any excursion; 14-day cycles otherwise & the inter-dose recovery is the only excursion a fixed schedule supplies, and six-monthly visits do not resolve it \\
excursion, required & $m$ to reach 14.6$\,\bar m$ from the operating point, or a band of fold range 9.1 straddling $K_m$ & the excitation condition. Below it the band is degenerate in the lumped ratio; above $K_m$ it is degenerate in $\theta$ against the load \\
excursion, available & inter-dose recovery reaches a fold of about 2.6, downward & so the required excursion is not available from the schedule alone \\
load record & protein intake to $\le 5\%$, dated, with enteral and parenteral sources separated & the load's unrecorded component is what makes saturation degenerate, and it is the noise term a marker panel cannot average away \\
labelled-propionate test & once per patient at entry, repeated at any dose change & the only measurement that anchors capacity independently of level. Reference precision 12.8\% in patients \\
assay precision & coefficient of variation $\le$ 5\%, reported on the log scale & a constant CV is a constant variance on the log scale, which is the only scale on which $\sigma$ has unambiguous units \\
instrument events & $n \ge 77$ per patient for $20\%$ on $\theta$ & the floor is set by $\sigma$ alone and no excursion beats it \\
cohort & $N \ge 3,121$ to place the fraction with $\mu<0$ within $\pm 10$ points, at the accessible design's sign error of 41\% near the boundary & sized on the sign error of $\mu$, not on a confidence interval for $\theta$. Per-patient misclassification attenuates the fraction towards one half and inflates the corrected variance by $(1-2p)^{-2}$, which diverges as $p \to 1/2$ \\
the asymmetry & the two errors are not equal and the costly one is larger & at equal $|\mu|$ either side of the boundary the simulated error is $6\%$ for treating a patient who did not need it and $53\%$ for failing to treat one who did. The design fails towards declaring patients safe \\
\bottomrule
\end{tabular}
% the terminal specification; every row is generated
\end{table}
\end{adjustwidth}

\subsection{Tier 1: the measurement the method's feasibility turns on}
\label{sec:tier1}

\textbf{What to run.} One cohort, two $1$-$^{13}$C-propionate breath tests bracketing a
single dose, on one day. The quantity measured is the repeatability of the assay
\emph{within a day and within a state}, which is what a pre/post pair around a dose depends
on.

\textbf{Why it is first.} By Section~\ref{sec:excitation} the therapy channel is what makes
$\theta$ identifiable, and \eqref{eq:dosereq} is stated in units of $\theta$: the ratio
$\Delta/\theta$ is only as well known as $\Delta$ is. Section~\ref{sec:designsec} turns
that into one number --- the delivered increment must be characterised to about $12\%$ for
monthly sampling to reach the $20\%$ target on $\theta$. \textbf{Everything the therapy
channel buys is conditional on that, and it has never been measured.}

\textbf{What is published is an upper bound, and it is not the quantity needed.} The only
repeatability data for this assay~\cite{breath_test_mma} % C48
are between-visit: controls repeated at one- to three-week intervals, and fifteen
\emph{mut} patients at separate visits \emph{two months to four years apart}, with the
authors naming changing renal status and varying $V\mathrm{CO}_2$ as causes. The Results
give $14.02 \pm 12.86\%$ with a range of $0.39$ to $50.84\%$; the abstract's widely quoted
$12.8\%$ is the standard deviation of the per-subject coefficients, not their mean. And the
cohort is methylmalonic acidemia, not propionic acidemia, where the block sits at PCC
itself.

Between-visit variation is an \emph{upper bound} on same-day error, because it contains
drift a same-day pair removes, and the bound sits close to the requirement. \textbf{But a
bound that happens to sit near a threshold is not a measurement of it}, and a per-subject
range reaching $50.8\%$ cannot be applied to an individual. Tier~1 replaces an inference
with a number.

\subsection{Tier 2: conditional on Tier 1, and then it is a small study}

\paragraph{Cover the dose cycle; do not chase the marker.} Monthly draws across dose cycles
for two years reach the target, and two-monthly draws --- thirteen over two years --- also
reach it. This is ordinary outpatient blood work. What buys the information is the
\emph{phase}, not the rate: a schedule that sees the increment vary identifies $\theta$ and
one that does not, does not, whatever its density.

\paragraph{And one hard exclusion a protocol would currently get wrong.}
\textbf{Do not sample on dosing days, or at any multiple of the dose interval}
(Remark~\ref{rem:aliasing}). ``Draw bloods when the patient attends for dosing'' is the
natural thing to write and it is the one schedule that destroys the result: every draw
lands at the same phase, the range of $\Delta/\theta$ collapses from $1.95$ to $0.13$, and
$\operatorname{sd}(\ln\hat\theta)$ goes from $0.17$ to $3.51$. Moving the schedule half a
day recovers it.

\paragraph{The load record: ordinary precision, and keep the variation.} Recording intake
twenty times more precisely buys about six per cent, so precision is not the binding
variable. The variability is subtler and enters twice. By Proposition~\ref{prop:rank} a
perfectly steady load makes the design exactly rank deficient, and that null direction must
be filled by something; as the specification stands it is filled by the $5\%$ prior on the
load bias, and only real day-to-day variation lets the \emph{data} fill it instead. So
``record intake as eaten, do not smooth it'' is the difference between a result the data
support and one a prior supports. It costs nothing: it is an instruction to omit a
data-cleaning step, not to add a measurement.

\paragraph{$k_e$, and nothing else, is assumed.} The result needs the non-enzymatic
clearance within a factor of about three, which the literature supplies; no measurement
programme is proposed for it. $K_m$ is assumed at nothing: a prior twenty times looser
leaves every number unchanged, which is what Remark~\ref{rem:no_km} demanded.

\subsection{What is excluded, and what is not promised}

Two acquisitions an earlier draft asked for are struck out. \textbf{High-frequency sampling
at a fixed operating point} buys nothing: two-monthly sampling with thirteen draws beats
fortnightly sampling with fifty-three, because the first sees the dose cycle and the second
is aliased against it. \textbf{Decompensation records} are excluded twice over --- the
excursion is not observed at the magnitude required, and by Section~\ref{sec:discussion}
the model does not hold in that regime.

\textbf{And the specification does not promise a decision.} It buys an identified $\theta$.
Section~\ref{sec:decision} shows that the sign of the margin does not follow from it, for
reasons that are not about the data, and Section~\ref{sec:scope} shows that $\theta$ itself
stops being recoverable in the patients the question matters most for. A protocol written
from this page should report $\hat\theta$ with its interval and should not be read as a
rule for escalating an individual patient.

\vspace{0.5em}
\hrule
\vspace{1em}

%==========================================================================
